\scanpage{143}
i.e. qu'on ait $f\delta x=\delta fx$, i.e. compte tenu de (4),
\[
d_E^1(fx)+(-fd_E^0x-x\otimes df)=d_E^1(fx)-d_E^0(fx),
\]
ce qui s'écrit
\[
d_E^0(fx)=fd_E^0(x)+x\otimes df \tag*{(7)}
\]
ou aussi, utilisant la structure de $\Omega^{\unsure{\cdot|\bullet}}_{S/T}$-module
à droite de $E\otimes\Omega^{\cdot}_{S/T}$ :
\[
d_E^0(\deleted{fx}xf)=d_E^0(x)f+x\,df.\tag*{(\deleted{\unsure{7}}7')}
\]

NB On dit souvent que $d_E^0$ (5) satisfaisant (7), est une connex. sur $E$ rel. à $T$.
\deleted{On sait alors que la \unsure{que}}

On sait alors qu'il $\exists$ prolongement unique de
$d_E^0$ en \deleted{une \unsure{dériv. extérieure|dérivation extérieure} \unsure{hom}} un endomorphisme
additif
\[
d_E^{\cdot}:E\otimes\Omega^{\cdot}_{S/T}\longrightarrow E\otimes\Omega^{\cdot}_{S/T} \tag*{(8)}
\]
qui \deleted{\unsure{est une anti}} satisfait à la relation suivante
rel. à la structure de $\Omega^{\cdot}_{S/T}$-module à droite
\[
d_E^{\cdot}(\varphi\omega)=d_E^{\cdot}(\varphi)\omega+
 (-1)^{\deg\varphi}\varphi\,d(\omega), \tag*{(9)}
\]
et ce $d_E^{\cdot}$ est un op de degré $1$. Pour que
l'on ait
\[
d_E^{\cdot}\circ d_E^{\cdot}=0, \tag*{(10)}
\]
il f. et s. que l'on ait
\[
\text{le composé }\ d_E^1d_E^0:E\longrightarrow E\otimes\Omega^1_{\unsure{S|E}/T}
\longrightarrow E\otimes\Omega^{\unsure{2}}_{\unsure{E|S}/T}\ \text{est nul.} \tag*{(10')}
\]

Le composé en question, considéré comme section
de $\operatorname{Hom}(E,E\otimes\Omega^2_{S/T})$, est appelé le \uline{tenseur courbure}
de la connexion. La connexion est dite \uline{intégrable}
si \unsure{le|la} \deleted{tenseur} courbure est \unsure{nul|nulle}, i.e. si \inserted{$d_E^{\cdot}$} (8) fait
% archive sheet number: 142
